\documentclass[11pt]{article}
\usepackage{amsmath,amssymb}
\usepackage[margin=1in]{geometry}
\title{Vertical semi-implicit acoustic solve on the RiRk (cubic B-spline) grid}
\author{Scythe.jl --- implementation note}
\date{2026-06-19}

\begin{document}
\maketitle

\section{Why the Chebyshev solver does not transfer directly}

The semi-implicit acoustic step (\texttt{semiimplicit.jl}) reduces, in each
vertical column, to a one-dimensional Helmholtz boundary-value problem
\begin{equation}
  \mathcal{L}\,\phi \;\equiv\; \bigl(\alpha\,\partial_{zz} + \beta\bigr)\phi \;=\; g(z),
  \qquad z\in[z_b,z_t],
  \label{eq:helmholtz}
\end{equation}
with one boundary condition at each of $z_b$ (bottom) and $z_t$ (top). The
unknown is represented by a vertical spectral coefficient vector $\mathbf{a}$,
and the field is recovered at the physical (mish) points by an inverse transform.

On the \textbf{RZ grid} the vertical basis is Chebyshev and the inverse transform
(the DCT matrix) is \emph{square}: the number of collocation points equals the
number of coefficients, $k_{\mathrm{Dim}} = b_{k\mathrm{Dim}}$, and the
coefficient vector $\mathbf{a}$ also has length $k_{\mathrm{Dim}}$. Writing
$M_0,M_1,M_2$ for the matrices that map coefficients to values / first / second
derivatives at the collocation points (each $k_{\mathrm{Dim}}\times
k_{\mathrm{Dim}}$), \eqref{eq:helmholtz} discretises to the square system
\begin{equation}
  H = \alpha M_2 + \beta M_0,
  \qquad
  H_a =
  \begin{bmatrix} \mathrm{BC}_b \\ \mathrm{BC}_t \\ H[\,2{:}n_z\!-\!1,:\,] \end{bmatrix},
  \qquad
  \mathbf{a} = H_a^{-1}
  \begin{bmatrix} 0 \\ 0 \\ g[\,2{:}n_z\!-\!1\,] \end{bmatrix},
  \label{eq:cheb}
\end{equation}
i.e.\ the two boundary rows of the collocated operator are replaced by the
boundary-condition rows, and the interior mish values of $g$ are the RHS.

The \textbf{RiRk grid} uses a cubic B-spline vertically, where
\begin{equation}
  b_{k\mathrm{Dim}} = n_c + 3, \qquad k_{\mathrm{Dim}} = n_c\,\bar{\mu},
\end{equation}
with $n_c$ cells and $\bar\mu$ Gauss (mish) points per cell. Now the basis
matrices $M_0,M_1,M_2$ from \texttt{operator\_matrix(:k,$\cdot$)} are
\emph{rectangular} $k_{\mathrm{Dim}}\times b_{k\mathrm{Dim}}$ (they evaluate the
$b_{k\mathrm{Dim}}$ basis functions at the $k_{\mathrm{Dim}}$ mish points), with
$k_{\mathrm{Dim}}>b_{k\mathrm{Dim}}$. Applying \eqref{eq:cheb} verbatim makes
$H_a$ rectangular and ``$H_a\backslash g$'' silently becomes an over-determined
least-squares solve, which does not satisfy \eqref{eq:helmholtz} and blows up.

We need a \emph{square} $b_{k\mathrm{Dim}}\times b_{k\mathrm{Dim}}$ system in
spline-coefficient space. Two routes were considered.

\section{Route 1: Greville (knot) collocation}

\paragraph{Greville abscissae.} For a B-spline of order $p+1$ (cubic: $p=3$) on
the knots $\{t_i\}$, the \emph{Greville abscissa} of the $i$-th basis function is
the average of its $p$ interior knots,
\begin{equation}
  \xi_i = \frac{1}{p}\sum_{j=1}^{p} t_{i+j}.
  \label{eq:greville}
\end{equation}
These are the classical collocation points for spline interpolation and
collocation BVP solvers: collocating a degree-$p$ spline at its
$b_{k\mathrm{Dim}}$ Greville abscissae gives a square, banded, well-conditioned
system (de~Boor). For the \emph{uniform} knots used here ($t_j=z_b+j\,\Delta z$,
basis function $m$ centred on knot $m$, $m=-1,\dots,n_c+1$), \eqref{eq:greville}
collapses to the central knot, $\xi_m = z_b + m\,\Delta z$. Two of these
($m=-1$ and $m=n_c+1$) fall outside $[z_b,z_t]$; we keep the $n_c+1$ in-domain
ones, which coincide with the cell-boundary nodes $\zeta_c = z_b+(c-1)\Delta z$,
and collocate the operator there.

\paragraph{System.} With $N_r=[\,\partial_z^{\,r}B_k(\zeta_c)\,]$ the
$(n_c+1)\times b_{k\mathrm{Dim}}$ node matrices,
\begin{equation}
  H_a=\begin{bmatrix}\mathrm{BC}_b\\\mathrm{BC}_t\\ \alpha N_2+\beta N_0\end{bmatrix}
  \in\mathbb{R}^{(n_c+3)\times(n_c+3)} ,
\end{equation}
square because $2+(n_c+1)=n_c+3=b_{k\mathrm{Dim}}$. The RHS, known at the mish
points, is resampled to the nodes with a natural cubic spline.

\paragraph{Verdict.} The static solve is exact (manufactured-solution error
$\sim10^{-6}$) and well-conditioned ($\kappa\sim10$--$10^2$). \emph{But the time
integration is unstable}: a slow acoustic mode grows and the run NaNs after
$\sim100$ steps. The node-collocation operator is non-symmetric and --- more
importantly --- it is built at the \emph{nodes}, whereas the model's explicit
acoustic tendencies use the spline derivative evaluated at the \emph{mish}
points (\texttt{gridTransform}). The AI2$^\ast$ implicit/explicit split requires
these two operators to coincide (they do for Chebyshev: one DCT serves both).

\section{Route 2: Galerkin (implemented)}

The fix is a Galerkin (finite-element) discretisation built from the
\emph{same} mish operators the explicit tendencies use. The weak form of
\eqref{eq:helmholtz}, tested against the spline basis and integrated with the
spline's own Gauss quadrature (physical weights $W$ at the mish points), is
\begin{equation}
  \underbrace{\bigl(-\alpha\,M_1^{\mathsf T}W M_1 + \beta\,M_0^{\mathsf T}W M_0\bigr)}_{A}\,
  \mathbf{a}
  \;=\;
  \underbrace{M_0^{\mathsf T} W\,g_{\text{mish}}}_{b},
  \label{eq:galerkin}
\end{equation}
using $\int\psi\,\partial_{zz}\phi = -\int\psi'\phi'$ (the boundary term vanishes
for the homogeneous conditions here). $A$ is square ($b_{k\mathrm{Dim}}\times
b_{k\mathrm{Dim}}$) and \emph{symmetric}; $M_0,M_1$ are exactly the value and
derivative operators \texttt{gridTransform} applies, so the implicit and explicit
acoustic operators are consistent. Neumann conditions are natural (no row
replacement); Dirichlet conditions replace the first/last rows of $A$ with the
boundary-value constraint $B_k(z_b)\mathbf a=0$, $B_k(z_t)\mathbf a=0$, and zero
the matching entries of $b$.

The three solver matrices map to \eqref{eq:galerkin} as
\[
  \text{(xi-form)}\ \alpha=-\tau^2\bar P_\xi,\ \beta=1;\quad
  \text{(w-form)}\ \alpha=\tau^2\bar P_\xi,\ \beta=-1;\quad
  \text{(diffusion)}\ \alpha=-\tau,\ \beta=1,
\]
with $\tau$ the time-stepping coefficient and $\bar P_\xi$ the domain-mean sound
speed squared. Manufactured-solution errors are $\sim10^{-6}$ (Dirichlet and
Neumann), and the operator is symmetric to machine precision.

\paragraph{Timestep.} With the Galerkin solver the RiRk benchmarks integrate
stably and reproduce the physical evolution, but the cubic-B-spline acoustic
solve has a \emph{tighter} stability limit than the Chebyshev pseudospectral
solve at the same horizontal resolution: at $200$~m cells the dry thermal is
unstable at $\Delta t=0.2$~s (the Chebyshev value) and stable at $\Delta
t=0.1$~s. The implicit vertical solve acts in the $b_{k\mathrm{Dim}}$ coefficient
space rather than collocating the field point-for-point, so the vertical acoustic
modes are not as completely removed from the explicit horizontal step as in the
pseudospectral case. The benchmark harness therefore halves $\Delta t$ for the
RiRk grid (\texttt{vertical\_ts}); physical times and outputs are unchanged.
Establishing the exact stability limit (and whether a fully consistent-mass or
mixed formulation recovers the Chebyshev $\Delta t$) is left as follow-up.

\section{Code map}

\begin{itemize}
  \item \texttt{\_rirk\_solve\_data} --- caches $M_0,M_1$ (mish operators), the
        quadrature weights $W$, and the boundary-value rows.
  \item \texttt{\_assemble\_vertical\_matrix} --- Chebyshev collocation
        (unchanged) or the spline Galerkin assembly \eqref{eq:galerkin}.
  \item \texttt{\_vertical\_solve!} --- forms the RHS ($g[2{:}n_z\!-\!1]$ for
        Chebyshev, $M_0^{\mathsf T}Wg_{\text{mish}}$ for the spline, with
        Dirichlet rows zeroed) and solves for $\mathbf a$.
\end{itemize}
All acoustic ($\xi$/$w$) and implicit-diffusion ($s,\mu,\dots$) solves route
through these helpers, so one code path serves both vertical bases.

\end{document}
